\documentclass[11pt]{article}
\usepackage[a4paper,margin=1in]{geometry}
\usepackage{amsmath,amssymb,mathtools,bm}
\usepackage{enumitem,booktabs,array}
\usepackage{tcolorbox}
\usepackage{hyperref}
\usepackage[protrusion=true,expansion=false]{microtype}
\hypersetup{colorlinks=true,linkcolor=blue,urlcolor=blue,citecolor=blue}
\setlist[itemize]{topsep=2pt,itemsep=2pt}
\setlist[enumerate]{topsep=2pt,itemsep=2pt}
\newtcolorbox{keybox}{colback=blue!3!white,colframe=blue!40!black,boxrule=0.4pt,arc=1mm}
\newtcolorbox{alertbox}{colback=red!3!white,colframe=red!40!black,boxrule=0.4pt,arc=1mm}
\newtcolorbox{answerbox}{colback=green!3!white,colframe=green!40!black,boxrule=0.4pt,arc=1mm}
\newtcolorbox{drillbox}{colback=yellow!5!white,colframe=orange!60!black,boxrule=0.4pt,arc=1mm}
\newcommand{\EV}{\mathbb{E}}
\newcommand{\Var}{\mathrm{Var}}
\newcommand{\Cov}{\mathrm{Cov}}
\newcommand{\blank}[1]{\underline{\hspace{#1}}}

\title{E10-01: Gompers, Kovner, Lerner \& Scharfstein (2010, JFE) \\
\large ``Performance Persistence in Entrepreneurship'' --- Active Recall Workbook}
\author{Empirical Corporate Finance II --- Exam Prep}
\date{\today}
\begin{document}\maketitle

\begin{keybox}
\textbf{Paper in one sentence.} GKLS show that entrepreneurs with prior successes (IPO or high-value acquisition) have $\sim$30\% probability of success in their next venture vs.\ $\sim$21\% for first-time entrepreneurs and $\sim$22\% for failed-prior entrepreneurs --- striking persistence in entrepreneurial performance driven by both inherent skill and network/reputation effects.

\textbf{Placement.} Canonical paper on entrepreneurial serial success; anchor for studies of founder experience, VC screening, and entrepreneurial human capital. Influences Kaplan-Sensoy-Str\"omberg (2009) on VC evaluation.
\end{keybox}

\section*{Round 1: Complete Walk-Through}

\subsection*{1.1 Research question}
Is entrepreneurial success persistent across ventures? If so, is persistence due to skill, industry/market timing, or VC-network effects? Answers matter for VC screening, public-policy support of entrepreneurship, and human-capital theory.

\subsection*{1.2 Data}
\begin{itemize}
\item VentureOne / Thomson Reuters VC-backed US startup data, 1986--2003.
\item Founder-level records linked across ventures.
\item Outcome: IPO or profitable acquisition (exit event), $\sim$20\% base rate.
\end{itemize}

\subsection*{1.3 Empirical design}
\begin{itemize}
\item Classify founders by prior-outcome: first-time, prior failure, prior success.
\item Logit regression of current-venture success on prior-outcome plus controls (industry, year, VC-backer).
\item Decompose persistence into skill vs.\ industry/market-timing components using industry fixed effects.
\end{itemize}

\subsection*{1.4 Headline results}
\begin{itemize}
\item Success rates: first-time 21\%, prior-failed 22\%, prior-succeeded 30\%.
\item Persistence is economically large and statistically significant.
\item Successful founders return to similar industries (skill + industry knowledge).
\item Part of persistence driven by market timing: successful founders cluster in high-performing industries / time periods.
\item Top-VC backing amplifies persistence (network + resources interact with skill).
\end{itemize}

\subsection*{1.5 Interpretation}
\begin{itemize}
\item Serial success cannot be fully explained by luck.
\item Mix of inherent skill, industry knowledge, and network access.
\item Challenges the view that entrepreneurship outcomes are dominated by randomness.
\end{itemize}

\subsection*{1.6 Implications}
\begin{itemize}
\item VC screening: founder track record is a valuable signal.
\item Policy: supporting failed entrepreneurs does not show a persistent penalty; failure is not highly informative of future performance.
\item Human-capital and social-capital both matter in entrepreneurial production functions.
\end{itemize}

\subsection*{1.7 Caveats}
\begin{alertbox}
\begin{itemize}
\item Sample limited to VC-backed companies (selected sample).
\item Success measured by IPO/acquisition, not long-run profitability.
\item Cannot fully distinguish skill from VC-network effects.
\end{itemize}
\end{alertbox}

\section*{Round 2: Level 1 Fill-in}
\begin{enumerate}
\item First-time success rate: \blank{1cm}\%.
\item Prior-failed success rate: \blank{1cm}\%.
\item Prior-succeeded success rate: \blank{1cm}\%.
\item Part of persistence: \blank{2cm} timing.
\item Sample: VC-backed US startups \blank{1cm}--\blank{1cm}.
\end{enumerate}

\section*{Round 3: Level 2 Prompts}
\begin{enumerate}
\item State the research question and success measure.
\item Describe the classification of founders.
\item Report persistence estimates.
\item Decompose into skill, timing, networks.
\item Discuss implications for VC screening and policy.
\end{enumerate}

\section*{Round 4: Minimal Scaffolding}
From a blank page: (i) question, (ii) data, (iii) persistence, (iv) mechanism, (v) implications.

\section*{Round 5: Blank Page}
Essay: entrepreneurial skill, persistence, and VC screening.

\vspace{6pt}
\begin{drillbox}
\textbf{Speed Drill.} (1) Success measure. (2) First-time rate. (3) Prior-success rate. (4) Key mechanism. (5) Policy implication.
\end{drillbox}

\section*{Quick Reference \& Answer Key}
\begin{answerbox}
\begin{itemize}
\item \textbf{Rates.} First-time 21\%, failed 22\%, success 30\%.
\item \textbf{Gap.} $\sim$9 pp persistence.
\item \textbf{Drivers.} Skill + timing + VC networks.
\item \textbf{Lesson.} Track record matters for screening.
\end{itemize}
\end{answerbox}

\begin{keybox}
\textbf{30-second exam pitch.} Gompers, Kovner, Lerner \& Scharfstein (2010) document striking performance persistence in entrepreneurship using VentureOne data on VC-backed US startups 1986--2003. Linking founders across ventures, they find that previously-successful founders (prior IPO or profitable acquisition) have $\sim$30\% success probability on their next venture, versus $\sim$21\% for first-time founders and $\sim$22\% for previously-failed founders. The $\sim$9 pp persistence gap is economically and statistically large, cannot be fully explained by industry/market timing, and reflects a combination of inherent skill, industry knowledge, and VC-network access. The paper anchors entrepreneurial-human-capital research, provides a clear rationale for VC screening on founder track record, and notes that prior failure does not show a measurable penalty --- relevant for policy toward second-chance entrepreneurship.
\end{keybox}

\end{document}
